\documentclass[varwidth=16cm, border=2mm]{standalone}
\usepackage{amsmath}\usepackage{amssymb}\usepackage{ctex}
\usepackage{tasks}\usepackage{xcolor}\usepackage{graphicx}
\settasks{label=\Alph*., label-width=1.5em, item-indent=2em}
\begin{document}
\textbf{[0.6]} （2024·全国·高三对口高考）已知椭圆$\frac{x^2}{9} + y^2 = 1$，过左焦点$F$作倾斜角为$\frac{\pi}{6}$的直线交椭圆于$A$、$B$两点，则弦$AB$的长为$_{\_}$.

\vspace{0.2cm}\par\noindent\textcolor{red}{\textbf{【答案】} 2}
\par\noindent\textcolor{blue}{\textbf{【解析】} 1. **确定椭圆参数和焦点坐标**
   椭圆方程为$\frac{x^2}{9} + y^2 = 1$。
   比较标准方程$\frac{x^2}{a^2} + \frac{y^2}{b^2} = 1$，得到$a^2 = 9$，$b^2 = 1$。
   所以$a = 3$，$b = 1$。
   计算焦距$c$：$c^2 = a^2 - b^2 = 9 - 1 = 8$，所以$c = \sqrt{8} = 2\sqrt{2}$。
   左焦点$F$的坐标为$(-c, 0) = (-2\sqrt{2}, 0)$。

2. **确定直线方程**
   直线过左焦点$F(-2\sqrt{2}, 0)$，倾斜角为$\frac{\pi}{6}$。
   直线的斜率$k = \tan(\frac{\pi}{6}) = \frac{\sqrt{3}}{3}$。
   直线的点斜式方程为$y - 0 = \frac{\sqrt{3}}{3}(x - (-2\sqrt{2}))$，即$y = \frac{\sqrt{3}}{3}(x + 2\sqrt{2})$。

3. **联立直线与椭圆方程，求解弦长**
   将直线方程代入椭圆方程$\frac{x^2}{9} + y^2 = 1$：
   $\frac{x^2}{9} + \left(\frac{\sqrt{3}}{3}(x + 2\sqrt{2})\right)^2 = 1$
   $\frac{x^2}{9} + \frac{3}{9}(x + 2\sqrt{2})^2 = 1$
   $\frac{x^2}{9} + \frac{1}{3}(x^2 + 4\sqrt{2}x + 8) = 1$
   两边同乘以$9$：
   $x^2 + 3(x^2 + 4\sqrt{2}x + 8) = 9$
   $x^2 + 3x^2 + 12\sqrt{2}x + 24 = 9$
   $4x^2 + 12\sqrt{2}x + 15 = 0$
   设直线与椭圆的两个交点为$A(x_1, y_1)$和$B(x_2, y_2)$。
   根据韦达定理，有：
   $x_1 + x_2 = -\frac{12\sqrt{2}}{4} = -3\sqrt{2}$
   $x_1 x_2 = \frac{15}{4}$

   弦长$AB$的计算公式为$AB = \sqrt{(x_2 - x_1)^2 + (y_2 - y_1)^2}$。
   由于$y_2 - y_1 = k(x_2 - x_1)$，所以$AB = \sqrt{(x_2 - x_1)^2 + k^2(x_2 - x_1)^2} = \sqrt{1+k^2}|x_2 - x_1|$。
   首先计算$|x_2 - x_1|$：
   $|x_2 - x_1| = \sqrt{(x_1 + x_2)^2 - 4x_1 x_2} = \sqrt{(-3\sqrt{2})^2 - 4 \cdot \frac{15}{4}}$
   $= \sqrt{18 - 15} = \sqrt{3}$
   然后计算$1+k^2$：
   $1+k^2 = 1 + \left(\frac{\sqrt{3}}{3}\right)^2 = 1 + \frac{3}{9} = 1 + \frac{1}{3} = \frac{4}{3}$
   将这些值代入弦长公式：
   $AB = \sqrt{\frac{4}{3}} \cdot \sqrt{3} = \frac{2}{\sqrt{3}} \cdot \sqrt{3} = 2$

   **方法二：使用焦点弦长公式**
   对于椭圆$\frac{x^2}{a^2} + \frac{y^2}{b^2} = 1$，通过焦点的弦长公式为$L = \frac{2ab^2(1+k^2)}{b^2 + a^2k^2}$，其中$k$是焦点弦的斜率。
   本题中，$a=3, b=1, k=\frac{\sqrt{3}}{3}$。
   $L = \frac{2 \cdot 3 \cdot 1^2 \left(1+\left(\frac{\sqrt{3}}{3}\right)^2\right)}{1^2 + 3^2 \left(\frac{\sqrt{3}}{3}\right)^2}$
   $L = \frac{6 \left(1+\frac{3}{9}\right)}{1 + 9 \cdot \frac{3}{9}}$
   $L = \frac{6 \left(1+\frac{1}{3}\right)}{1+3}$
   $L = \frac{6 \cdot \frac{4}{3}}{4}$
   $L = \frac{8}{4} = 2$

   两种方法计算结果一致，弦$AB$的长为$2$。}


\end{document}
